\documentclass[aspectratio=169,10pt,english]{beamer}

\input{../../../_preambulo_beamer.tex}
\input{../../../_comandos_beamer.tex}

\selectlanguage{english}

\title{MA1001B - Statistical Modeling for Decision Making}
\subtitle{Lesson 38: Paired Samples (Before-After t)}
\author{}
\institute{Tecnol\'ogico de Monterrey}
\date{}

\begin{document}

\begin{frame}[plain]
  \vspace{1.2cm}
  {\Large\bfseries MA1001B - Statistical Modeling for Decision Making\par}
  \vspace{0.7cm}
  {\large Lesson 38: Paired Samples (Before-After t)\par}
  \vspace{0.7cm}
  {\color{TecRojo}\rule{\textwidth}{1.2pt}}
  \vspace{0.8cm}

  MA1001B Faculty\\[0.3cm]
  Term\\[0.3cm]
  Tecnol\'ogico de Monterrey
\end{frame}

\begin{frame}{The Before-After Question}
  \begin{alertblock}{Guiding question}
    How do we test a layout change when the same 25 stations are measured
    before and after?
  \end{alertblock}

  \vspace{0.6em}
  By the end of this lesson, you can:
  \begin{itemize}
    \item form paired differences $d=\text{before}-\text{after}$;
    \item run a one-sample $t$ test of $H_0:\mu_d=0$;
    \item contrast the paired SE with an unpaired two-sample SE;
    \item explain why unpaired analysis ignores the matching.
  \end{itemize}

  \vfill
  \scriptsize All handle times used today are synthetic. No real company data are used.
\end{frame}

\begin{frame}{The Unit Is the Station, Not the Column}
  \small
  A synthetic packing cell records handle time on the same 25 stations
  before and after a layout change. The observational unit is the station.

  \[
    d_i = \text{before}_i - \text{after}_i,\qquad i=1,\ldots,25.
  \]

  \begin{exampleblock}{Paired hypotheses}
    $H_0:\mu_d=0$ versus $H_1:\mu_d\neq 0$, with $\alpha=0.05$. This is a
    one-sample $t$ on the 25 differences, not a two-sample $t$ on 50 times.
  \end{exampleblock}
\end{frame}

\begin{frame}[fragile]{Step 1: Matched Differences}
  \begin{columns}[T]
    \begin{column}{0.42\textwidth}
      \small
      \begin{lstlisting}
d = before - after
dbar = d.mean()
s_d = d.std(ddof=1)
      \end{lstlisting}

      \begin{block}{Verified output}
        $\bar{d}=2.720813$\\
        $s_d=1.525772$\\
        $r=0.973356$
      \end{block}
    \end{column}
    \begin{column}{0.58\textwidth}
      \centering
      \includegraphics[width=\textwidth]{../figures/paired_t_01_differences.png}
      \vspace{0.2em}
      \scriptsize Matched before-after lines from Step 1
    \end{column}
  \end{columns}
\end{frame}

\begin{frame}{Paired $t$ Is One-Sample $t$ on $d$}
  \small
  \[
    t=\frac{\bar{d}}{s_d/\sqrt{n}}=\frac{2.720813}{0.305154}=8.916186,
    \qquad df=24.
  \]
  \[
    p=2P(T_{24}\ge 8.916186)=4.395592\times 10^{-9}.
  \]
  Because $|t|>t^*=2.063899$ and $p<\alpha$, the paired test rejects
  $H_0:\mu_d=0$.
\end{frame}

\begin{frame}[fragile]{Step 2: Paired Decision}
  \begin{columns}[T]
    \begin{column}{0.40\textwidth}
      \small
      \begin{lstlisting}
se = s_d / sqrt(n)
t = dbar / se
p = 2 * t_dist.sf(abs(t), 24)
      \end{lstlisting}

      \begin{block}{Verified output}
        $SE=0.305154$, $t=8.916186$\\
        $p=4.395592\times 10^{-9}$\\
        Decision: reject $H_0$
      \end{block}
    \end{column}
    \begin{column}{0.60\textwidth}
      \centering
      \includegraphics[width=\textwidth]{../figures/paired_t_02_paired_test.png}
      \vspace{0.2em}
      \scriptsize Histogram of differences from Step 2
    \end{column}
  \end{columns}
\end{frame}

\begin{frame}{Step 3: Unpaired Analysis Ignores the Matching}
  \begin{columns}[T]
    \begin{column}{0.42\textwidth}
      \small
      Paired SE $=0.305154$, $p=4.395592\times 10^{-9}$, reject.

      \vspace{0.4em}
      Unpaired Welch SE $=1.847143$, $p=0.147291$, do not reject.

      \vspace{0.4em}
      Before-after correlation: $0.973356$.

      \vspace{0.5em}
      \textbf{Limit:} treating the two columns as independent samples
      discards the pairing and can flip the $\alpha=0.05$ decision.
    \end{column}
    \begin{column}{0.58\textwidth}
      \centering
      \includegraphics[width=\textwidth]{../figures/paired_t_03_unpaired_limit.png}
      \vspace{0.2em}
      \scriptsize Paired versus unpaired p-values from Step 3
    \end{column}
  \end{columns}
\end{frame}

\begin{frame}{Concept Check}
  \begin{enumerate}
    \item Why is $d_i=\text{before}_i-\text{after}_i$ the right observational
      unit here?
    \item Compute $SE=s_d/\sqrt{n}$ from $s_d=1.525772$ and $n=25$.
    \item If $p=4.395592\times 10^{-9}$ and $\alpha=0.05$, what is the paired
      decision?
    \item Why does the unpaired $p=0.147291$ not refute the paired finding?
  \end{enumerate}

  \vspace{0.6em}
  \begin{block}{Connection to Lesson 39}
    The session can continue directly into two independent proportions in an
    A/B test, where the samples are not paired.
  \end{block}

  \vfill
  \scriptsize Source: OpenStax, \textit{Introductory Business Statistics 2e},
  Chapter 10, Section 10.6. CC BY-NC-SA.
\end{frame}

\appendix



\end{document}
